\section{\bi{Counterexamples, conclusions and remaining boundary}{反例、結論及尚待確定的邊界}}
\bi{The hard floor selects exact contacts by a first-order variation; on the six-site strip the remaining freedom is a single relative displacement. A comparison of the two endpoints misses the entire intermediate phase. At $s=8$, $h=263/2000$ and $t=1/1000$, direct rational image sums give $E(Y(h,b))-E(Y(h,0))>23/100000$ and $E(Y(h,0))-E(Y(h,t))>10^{-8}$. At $h=3/20$ the order of the endpoints reverses: the five-contact trial beats the center by more than $.00129689$. At $s=8.5,h=.13$ set $g_i=h+(1-6h)u_i$. The off-face trial $2^{20}u=(0,0,184621,0,1,863954)$ beats the center by more than $.0000876382$, yet loses by more than $6.3349\,10^{-8}$ to the on-face trial $2^{40}u=(905905483609,0,0,193606144167,0,0)$. Each comparison is between explicit cells. The global conclusions require Theorem~\ref{th:strip}.}{硬下限藉一階變分選出精確接觸；在六粒子條帶上，餘下的自由度是一個相對位移。僅比較兩端構型會漏掉整個偏移相。在 $s=8$、$h=263/2000$、$t=1/1000$，有理像和給出 $E(Y(h,b))-E(Y(h,0))>23/100000$ 與 $E(Y(h,0))-E(Y(h,t))>10^{-8}$。在 $h=3/20$ 兩端的次序反轉，五接觸比居中者低逾 $.00129689$。在 $s=8.5,h=.13$，記 $g_i=h+(1-6h)u_i$，面外試探 $2^{20}u=(0,0,184621,0,1,863954)$ 比居中者低逾 $.0000876382$，卻比面內的 $2^{40}u=(905905483609,0,0,193606144167,0,0)$ 高逾 $6.3349\,10^{-8}$。此處是指定構型間的比較；全域結論使用定理~\ref{th:strip}。}

\bi{The collision competitor in \cite{infimum} beats the limiting dipole at width 22. At the positive floor $h=1/1000$, let $K$ be the feasible gap simplex and $F$ the closed packed-triple face with its cyclic copies. Define}{文獻 \cite{infimum} 的碰撞構型在寬度 22 勝過極限偶極。固定正下限 $h=1/1000$，令 $K$ 為可行間隙單純形，$F$ 為緊雙三簇閉面及其循環移位之聯集。定義}
\[
 D(s)=\min_{x\in F}E_s(x)-\min_{x\in K}E_s(x),\qquad
 s_*=\inf\{s>9:D(s)>0\}.
\]
\bi{Joint continuity of the energy and compactness of both fixed sets give continuity of $D$. Theorem~\ref{th:strip} gives $D(9)=0$, and the strict whole-face comparison in Theorem~\ref{th:point} gives $D(22)>0$. The set of widths with strict defeat is open. Its lower boundary $s_*$ remains unknown, as does coexistence there with an off-face minimizer. This definition allows more than one component. Theorem~\ref{th:point} classifies the contact type at the stated pair; an explicit surrounding phase region remains to be determined. For general $N$, Theorem~\ref{th:lock} gives a sufficient small-floor bound and leaves the endpoint of exact locking open.}{能量聯合連續，兩個固定集合皆緊，故 $D$ 連續。定理~\ref{th:strip} 給出 $D(9)=0$，定理~\ref{th:point} 的整面嚴格比較給出 $D(22)>0$。緊面嚴格落敗的寬度構成開集。其下邊界 $s_*$ 的數值，以及該處是否有面外極小者共存，仍待確定。此定義允許多個連通分支。定理~\ref{th:point} 在所述參數點確定接觸類型，周圍相區的顯式範圍仍待求出。對一般 $N$，定理~\ref{th:lock} 給出充分的小間距界，精確鎖定的終點仍未確定。}

\paragraph{\bi{Research and writing disclosure.}{研究及撰寫披露。}}
\bi{AI-assisted work produced trial configurations, symbolic derivations, proof drafts, interval programs and explanatory text at successive stages. This edition used Codex to integrate and translate those local materials, audit the mathematics, compile the articles and coordinate a separate agent review. Trial searches used floating-point arithmetic. Claims of strict inequality rely on analytic estimates or outward interval bounds. The appendix identifies the small formalized portion. The independent review in the archive is an AI review; journal and independent human refereeing remain separate steps.}{各研究階段使用 AI 輔助提出試探構型、推導公式、擬寫證明、編寫區間程序及說明文字。本版以 Codex 整合及翻譯本機底稿、複核數學、編譯論文，並安排獨立 Agent 審閱。試探搜索採浮點算術；嚴格不等式均以解析估計或向外捨入的區間界為依據。附錄逐項說明已形式化的部分。附件所載獨立審查由 AI 完成，期刊與獨立人類同行審查仍待進行。}

\appendix
\section{\bi{Certified bounds on the continuous configuration space}{連續構型空間的嚴格區間證書}}\label{sec:cert}
\bi{\emph{What a terminal box proves.} For $h<1/6$ write $S=1-6h$ and $g_i=h+Su_i$ with $u_i\ge0$, $\sum_i u_i=1$. Put a largest $u_i$ at position 5. The other positive indices form $A\subseteq\{0,\ldots,4\}$. Reflection has eight fixed subsets among 32, hence 20 representatives. The target closed edge and vertex have $A=\{2\}$ and $\varnothing$; 18 supports remain. For $C_k=\{(i,j):i\le k<j\}$, $C_5=\varnothing$, define $d_{ij}=(j-i)h+S\sum_{p=i}^{j-1}u_p$, $w_{ij}=(-1)^{i+j}$. The unordered score with kernel $F_s$ below must satisfy}{\emph{終止盒所證明的命題。} 對 $h<1/6$ 令 $S=1-6h$，$g_i=h+Su_i$、$u_i\ge0$、$\sum_i u_i=1$。將一個最大 $u_i$ 重標在第 5 號間隙，其餘正坐標構成 $A\subseteq\{0,\ldots,4\}$。32 個子集在反射下有 8 個不動者，故有 20 種代表。目標閉邊及頂點對應 $A=\{2\}$ 與空集，尚餘 18 種支撐。對 $C_k=\{(i,j):i\le k<j\}$、$C_5=\varnothing$，記 $d_{ij}=(j-i)h+S\sum_{p=i}^{j-1}u_p$、$w_{ij}=(-1)^{i+j}$。以下文 $F_s$ 作核，無序得分須滿足}
\begin{equation}\label{eq:cuts}
 \sum_{C_k} w_{ij} F_s'(d_{ij})=0,\qquad
 \sum_{C_k\triangle C_l}w_{ij}F_s''(d_{ij})\le0 ,
\end{equation}
\bi{for transfers involving slack gaps (5 to $k$ for the gradient, $k$ to $l$ for curvature). Dividing out $S$ or $S^2$ is legitimate in this range. Both signs of the transfer are feasible in the original simplex even if two largest gaps tie. The cuts alter only crossing pairs; the squared transfer velocity selects the symmetric difference in the curvature. Closed coordinate boxes contain the entire relative-open face. A boundary point where a gap ceases to be slack is tested on its own support. Coordinate sums are intersected with bounds from their complements and the total-sum constraint. A box terminates if empty, strictly beaten by a feasible trial, or violating either necessary condition. Otherwise bisect the longest free coordinate with both children sharing the midpoint. An unresolved box at the minimum step fails the verifier. Success gives a finite proof by induction on the actual trees.}{其中涉及轉移的間隙須鬆弛（一階由第 5 隙轉至 $k$；二階在 $k,l$ 之間）。此時可除以 $S$ 或 $S^2$。即使最長間隙並列，原單純形中仍可雙向變分。只有跨越切口的配對改變，速度平方給出二階式的對稱差。閉坐標盒覆蓋整個面的相對內部；若某鬆隙在邊界變為接觸，該點另由自己的實際支撐處理。坐標和的包圍與補集及總和約束的包圍取交集。空盒、被可行試探嚴格擊敗或違反一項必要條件均可終止，否則沿最長自由坐標對分，兩子盒均含中點。達最小步長仍未判定即驗證失敗。成功時對所生成的有限樹歸納即得覆蓋證明。}

\bi{\emph{Uniform kernel bounds.} On $[\ell,r]$ with midpoint $a$, replace $G_s$ by $F_s=c_a+(c_a q_a)/(c_s q_s)(G_s-c_s)$, $q_s=e^{-\pi^2/s}$, $c_s=\sqrt{\pi/s}$. The slope is positive and the constant in the score is shape-independent. The first harmonic is fixed. Let $D_n$ be the larger absolute difference of $q_s^{n^2-1}$ at either endpoint from its value at $a$. For $j=0,1,2$ use}{\emph{寬度一致的核界。} 在以 $a$ 為中點的 $[\ell,r]$，以 $F_s=c_a+(c_a q_a)/(c_s q_s)(G_s-c_s)$ 代替 $G_s$，其中 $q_s=e^{-\pi^2/s}$、$c_s=\sqrt{\pi/s}$。斜率為正，得分的常數變化與構型無關，第一諧波保持固定。令 $D_n$ 為兩端 $q_s^{n^2-1}$ 與中點值所差的較大絕對值。對 $j=0,1,2$ 用}
\[
 \|F_s^{(j)}-G_a^{(j)}\|_\infty\le
 2c_a q_a\left(\sum_{n=2}^4(2\pi n)^jD_n+
 \frac{(10\pi)^j q_r^{24}}{1-(6/5)^j q_r^{11}}\right).
\]
\bi{The bounds below are in units of $10^{-9}$. On the narrow band they follow instead by differentiating $q_s^{n^2-1}$ and summing the geometric tail relative to width 8. At the reference width tabulate images $|p|\le3$ at mesh $1/8192$. Neighboring Gaussian values are generated by products of ratios whose own ratio is constant. Enclose the starting exponentials by alternating Taylor sums and round each product outwards to units of $10^{-32}$. The derivative multipliers are $1,-2a d,(2ad)^2-2a$ at image argument $d$. Omitted images with $|p|\ge4$ on $[8,9]$ contribute at most $4(324\cdot25+18)2^{-72}<10^{-9}$ in each of the three orders. Spatial extensions use bounds $4,20,150$ on the first three absolute derivatives; at the reference width 8 the original fixed-width table has its own certified bounds. At a score midpoint the code uses the value and slope at the nearest node, a distance at most $1/(2M)$ away, $M=8192$. The first-order Taylor remainder is at most $\|G_a''\|_\infty/(8M^2)$. The node value and slope each have error at most $2\cdot10^{-9}$; flooring the product costs at most one more unit of $10^{-9}$. Widen both the box score and the trial score by $15\epsilon_0$ at the reference center, then bound variation by the widened derivatives. Trial allocations put $u_2=0,1/8,1/4,3/8,1/2$ with the remainder in $u_5$. Only a strict comparison prunes.}{下列各界均以 $10^{-9}$ 為單位。窄帶的界以寬度 8 為基準，對 $q_s^{n^2-1}$ 微分並累加幾何尾項而得。在參考寬度上以 $1/8192$ 網距製作 $|p|\le3$ 的像和表，相鄰高斯值依比值遞推，而比值自身成等比。起始指數用交替 Taylor 級數包圍，每次乘積向外捨入至 $10^{-32}$。對像位差 $d$，導數乘子依次為 $1,-2a d,(2ad)^2-2a$。在 $[8,9]$，被略去的 $|p|\ge4$ 像於三種階數各至多貢獻 $4(324\cdot25+18)2^{-72}<10^{-9}$。空間擴寬採前三階絕對導數界 $4,20,150$；參考寬度 8 的原始核表另有固定寬度認證界。在得分中點用最近節點的函數值及斜率作一階 Taylor 近似。取 $M=8192$，距離至多 $1/(2M)$，故餘項不逾 $\|G_a''\|_\infty/(8M^2)$。節點值及斜率的誤差各至多 $2\cdot10^{-9}$，乘積向下取整再預留一個 $10^{-9}$ 單位。參考中心的盒得分上界及試探下界各擴寬 $15\epsilon_0$，再用擴寬後的導數控制變化。試探取 $u_2=0,1/8,1/4,3/8,1/2$，餘量全在 $u_5$。僅嚴格比較可剪枝。}
\[
\begin{array}{r|rrr}
 [\ell,r]&10^9\epsilon_0&10^9\epsilon_1&10^9\epsilon_2\\\hline
 {[7.99,8.01]}&42761&539098&6806367\\
 {[8,8.5]}&1156027&14580271&184224460\\
 {[8.5,9]}&1309984&16545550&209498225
\end{array}
\]
\par
\begin{table}[tbp]\centering
\caption{\bi{Complete support trees. The last three columns count excluded leaves.}{完整支撐搜尋樹。後三欄為各類排除葉數。}}\label{tab:trees}
\begin{tabular}{lrrrrr}\toprule
$s$ & \bi{floor bands}{下限段} & \bi{boxes}{訪問盒} & \bi{score}{得分} & \bi{gradient}{梯度} & \bi{curvature}{曲率}\\\midrule
$[7.99,8.01]$ &14&6068&261&377&2522\\
$[8,8.5]$ &49&40740&3094&2938&14779\\
$[8.5,9]$ &49&39056&2844&2885&14240\\\bottomrule
\end{tabular}\end{table}
\par
\bi{The 14 narrow bands run from $.10$ to $.12$ in $.002$ steps, then to $.134,.14,.15,1/6$. For the broad bands $300000h$ runs from $21000$ to $49800$ in steps of $600$, then to $50000$. The leaf identity for the two broad trees is $2(5938+5823+29019)-2\cdot49\cdot18=79796$. Independent fixed-width full-simplex verifiers cover $h\ge1/50$ at 8 and $h\ge1/20$ at 9 up to the uniform endpoint. At 9 there are 60,182 visited boxes.}{窄帶的 14 段由 $.10$ 至 $.12$ 以 $.002$ 遞增，再至 $.134,.14,.15,1/6$。廣帶取 $300000h$ 自 $21000$ 至 $49800$ 每次增 $600$，末端再至 $50000$。兩棵廣帶樹的葉數恆等式為 $2(5938+5823+29019)-2\cdot49\cdot18=79796$。獨立固定寬度的全單純形證書於 8 覆蓋 $h\ge1/50$、於 9 覆蓋 $h\ge1/20$ 至均勻端點。寬度 9 的搜尋共訪問 60,182 盒。}

\bi{\emph{A separate rounding proof at width 8.} Round cumulative excesses $M\sum_{k<i}u_k$ down after one common uniform shift in $[0,1)$. All contacts are preserved. For each pair the physical error has mean zero, magnitude below $\rho=S/M$ and variance at most $\rho^2/4$ because a rounded pair difference exceeds its unshifted integer part by a Bernoulli variable whose success probability is the fractional part of the unrounded excess difference. Taylor's theorem for the fifteen unordered terms bounds the expected energy error by $(5/8)\|G''\|_\infty\rho^2$. Consequently some rounded cell has $E_{\rm round}\le E_s(x)+(5/8)\|G''\|_\infty\rho^2$. Finite rational comparisons over $[.08,.10]$ and $[.10,.132]$ confine a minimizer to within $\rho$ in each pair displacement of a viable packed-face grid cell. Transferring length across $m$ sites changes only $m(6-m)$ pair terms; the error in its energy derivative is at most $m(6-m)\|G''\|_\infty\rho/3$. The certified positive release derivatives force all four internal contacts. This second route illustrates why grid energy comparisons alone do not establish a continuum statement.}{\emph{寬度 8 的獨立取整證明。} 對累積餘量 $M\sum_{k<i}u_k$ 加同一個 $[0,1)$ 均勻隨機位移再向下取整，所有接觸均保留。取整後的累積餘量之差等於原差的整數部加上一個 Bernoulli 變量，其成功概率等於原差的小數部。因此各對物理位差誤差期望為零，絕對值小於 $\rho=S/M$，變異數至多 $\rho^2/4$。對十五項無序配對使用 Taylor 定理，期望能量誤差至多 $(5/8)\|G''\|_\infty\rho^2$，故至少一個取整構型滿足單側界 $E_{\rm round}\le E_s(x)+(5/8)\|G''\|_\infty\rho^2$。有理枚舉覆蓋 $[.08,.10]$ 及 $[.10,.132]$，將任一極小者逐對位差定位在某個未被排除的緊面格點的 $\rho$ 內。跨 $m$ 個粒子的轉移只改變 $m(6-m)$ 項，能量導數誤差至多 $m(6-m)\|G''\|_\infty\rho/3$。認證後放寬導數仍為正，迫使四個內間隙精確接觸。此獨立證明亦顯示，由格點能量比較推到連續構型還需要哪些條件。}

\bi{\emph{At width 22.} For blocks of lengths $a,b$ and initial signs $\epsilon_i,\epsilon_j$, form the combined kernel
\[K_{ab}(d)=\sum_{v=0}^{a-1}\sum_{w=0}^{b-1}(-1)^{v+w}G(d+(w-v)h).\] The test for a velocity vector $p$ of whole blocks sums $\epsilon_i\epsilon_j(p_j-p_i)K_{ab}'$ and $\epsilon_i\epsilon_j(p_j-p_i)^2K_{ab}''$ over block pairs. Add the fixed intrablock terms for energy tests. The first cover works in slack gaps of prescribed sum with adjacent-gap reflection symmetry; it checks nonempty proper subsets of blocks moving together. Its node spacing is $1/16000$, derivative intervals in units $10^{-9}$. The second enumerates all cyclic words, propagates $y_{j+1}\ge y_j+a_j h$ for block starts, includes both arc and complementary arc bounds, and also tests opposite block velocities. It uses mesh $1/32000$ and units $10^{-11}$. Retain seven images at each node, adding the two-sided tail $4(1936\cdot25+44)2^{-198}$ for orders zero through two. Compute the signed block sums before interpolating, enclosing the entire interval by widening node extrema by $\|K_{ab}^{(j+2)}\|_\infty/(8M^2)$. In the second cover this norm is bounded by $ab$ times $44,1500,22000$ for $j=0,1,2$. The two covers differ in coordinates, tables and cut directions. Their common analytic assumptions, including the image tail bound, remain open to inspection.}{\emph{寬度 22。} 對長度 $a,b$、起始符號 $\epsilon_i,\epsilon_j$ 的兩塊，合成塊核
\[K_{ab}(d)=\sum_{v=0}^{a-1}\sum_{w=0}^{b-1}(-1)^{v+w}G(d+(w-v)h).\]對整塊速度 $p$，將 $\epsilon_i\epsilon_j(p_j-p_i)K_{ab}'$ 與 $\epsilon_i\epsilon_j(p_j-p_i)^2K_{ab}''$ 分別對所有塊對求和。能量檢驗另加固定的塊內項。第一套覆蓋以總和固定的鬆隙為坐標，利用相鄰外隙反射對稱，檢驗一個非空真子集內的塊同時移動；其節點距 $1/16000$，導數區間單位為 $10^{-9}$。第二套枚舉全部循環字，在塊首坐標上傳播 $y_{j+1}\ge y_j+a_j h$，同時使用弧與補弧界，且另檢驗兩塊相反速度。網距為 $1/32000$，單位 $10^{-11}$。每個節點保留七個像，對零至二階補上雙向尾項 $4(1936\cdot25+44)2^{-198}$。先合成帶號塊和再插值，以 $\|K_{ab}^{(j+2)}\|_\infty/(8M^2)$ 擴寬節點極值，涵蓋整段實數位差。第二套對 $j=0,1,2$ 分別以 $ab$ 乘 $44,1500,22000$ 控制此導數範數。兩套覆蓋使用不同坐標、核表和變分方向；其共用的解析假設，包括像和尾界，均可獨立查核。}

\section{\bi{Formal coverage and reproduction}{形式化範圍與復算指引}}\label{sec:art}
\bi{The file \texttt{AlternatingMoments.lean} proves four declarations. The first two prove the two identities in \eqref{eq:mom}; the general inequality follows analytically by inserting the floor. The other two give the triple variance, its gap derivatives and their lower bounds. Their axiom lists contain exactly \texttt{propext}, \texttt{Classical.choice}, and \texttt{Quot.sound}. The base project's \texttt{Counterexample.lean} separately verifies the earlier four-site Gaussian counterexample. The global contact and phase theorems use the analytic proofs and, where specified, executable interval bounds.}{檔案 \texttt{AlternatingMoments.lean} 證明下列四項聲明。前兩項證明式 \eqref{eq:mom} 的兩個恆等式；一般不等式由代入間距下限解析推出。後兩項證明三粒子方差、兩個間隙導數及其下界。四項公理清單皆恰為 \texttt{propext}、\texttt{Classical.choice} 及 \texttt{Quot.sound}。基礎專案的 \texttt{Counterexample.lean} 另行完整驗證先前的四粒子高斯反例。全域接觸及相圖定理的依據則為上述解析證明及明列的可執行區間界。}

\begin{center}\small
\begin{tabular}{ll}
\texttt{firstMoment} & \texttt{secondMoment}\\
\texttt{tripleVariance\_eq} & \texttt{triple\_gap\_derivatives}
\end{tabular}
\end{center}

\bi{From the root of the public archive, run}{在公開附件的根目錄執行}
\begin{center}\small\texttt{python reproduce.py all -{}-work-dir ../chain-replay}\end{center}
\bi{The setup stage prepares the saved sources and dependencies in the specified work directory before the original verification programs run. The reproduction guide covers the interval certificates, numerical experiments, phase diagram, Lean proofs and both articles. The evidence index identifies the proof inputs for each theorem and distinguishes them from numerical comparisons. The citations and attribution in this edition refer to the source versions preserved in its local research record.}{設定階段先在指定工作目錄準備固定版本的來源及相依項，再執行原驗證程式。復算指南涵蓋區間證書、數值實驗、相圖、Lean 證明及兩份論文。證據索引逐一定理標明證明所用的輸入，並列出另作比較的數值實驗。本版引文與成果歸屬均以本次本地研究記錄所保存的來源版本為據。}

